\chapter{Appendix}

\section{Implementation of the ResNet Classifier}

This appendix presents the implementation of the ResNet-based classifier used in our experiments. The full implementation is provided in Listing~\ref{lst:resnet_code}. The code is written in Python using the PyTorch framework and follows the architectural requirements described in the project specification. 

The implemented model is a one-dimensional variant of the ResNet-18 architecture designed to process sequences of raw network packet bytes. Since network packets can be represented as sequences of byte values, a 1D convolutional architecture is used instead of the traditional 2D convolutional layers used in image processing.

The implementation consists of several main components:

\begin{itemize}

\item \textbf{Residual Blocks.}  
The model uses the BasicBlock structure of ResNet-18. Each block contains two convolutional layers followed by batch normalization and ReLU activation. A skip connection adds the input of the block to its output, enabling residual learning and improving gradient flow during training.

\item \textbf{ResNet Backbone.}  
The network contains four stages, each composed of two residual blocks, resulting in a total of 18 layers. The architecture follows the standard ResNet-18 layout adapted to 1D data. After the convolutional stages, an adaptive average pooling layer reduces the feature representation to a fixed-size embedding.

\item \textbf{Fully Connected Classification Head.}  
The extracted embedding is passed through two fully connected layers with dropout regularization before producing the final classification logits.

\item \textbf{Statistical Features.}  
In addition to raw packet bytes, the model optionally incorporates five statistical features derived from the packet content. These include:
\begin{itemize}
\item total number of bytes in the packet,
\item frequency of the most common byte,
\item frequency of the second most common byte,
\item Shannon entropy of the byte distribution,
\item fraction of zero bytes in the packet.
\end{itemize}

The entropy feature captures the randomness of the byte distribution, while the fraction of zero bytes helps detect padding patterns often introduced during packet normalization.

\item \textbf{Training Procedure.}  
The model is trained using the Adam optimizer and the categorical cross-entropy loss function. Early stopping is applied based on validation loss to prevent overfitting.

\item \textbf{Hyperparameter Search.}  
A grid search is performed over several training parameters, including the number of epochs, batch size, learning rate, and dropout probability. The combination yielding the highest average F1-score across the cross-validation folds is selected.

\item \textbf{Evaluation.}  
The implementation supports k-fold cross-validation and computes precision, recall, and F1-score for each fold. The final performance metrics are obtained by averaging the results across all folds.

\end{itemize}

The complete implementation used to run the ResNet experiments is provided below.

\begin{lstlisting}[language=Python, caption={Implementation of the ResNet classifier used in the experiments}, label={lst:resnet_code}]


import csv
from pathlib import Path

import numpy as np
import torch
import torch.nn.functional as F
from torch.utils.data import DataLoader, TensorDataset

from file_helper_t3 import load_k_fold_results
from labels_helper import deduplicate_folds, encode_labels
from handling_re_bytes_integrated import get_keep_indices_from_fold0_ae


# ---------------------------
# ResNet1D model definition
# ---------------------------

import torch
import torch.nn as nn
import torch.nn.functional as F


# -----------------------------
# 1) BasicBlock (1D) for ResNet-18
# -----------------------------
class BasicBlock1D(nn.Module):

    
    expansion = 1

    def __init__(self, in_channels: int, out_channels: int, stride: int = 1):
      
        # initialize parent class
        super().__init__()

      
        self.conv1 = nn.Conv1d(
            in_channels, out_channels,
            kernel_size=3, stride=stride, padding=1, bias=False
        )

        self.bn1 = nn.BatchNorm1d(out_channels)

        self.conv2 = nn.Conv1d(
            out_channels, out_channels,
            kernel_size=3, stride=1, padding=1, bias=False
        )

        # second batch normalization layer
        self.bn2 = nn.BatchNorm1d(out_channels)

        self.downsample = None
        if stride != 1 or in_channels != out_channels:
           
            self.downsample = nn.Sequential(
                nn.Conv1d(in_channels, out_channels, kernel_size=1, stride=stride, bias=False),
                nn.BatchNorm1d(out_channels),
            )

    def forward(self, x):
       
        # stores the input for skip connection
        identity = x

        # First conv + BN + ReLU
        out = self.conv1(x)
        out = self.bn1(out)
        out = F.relu(out, inplace=True)

        out = self.conv2(out)
        out = self.bn2(out)

        if self.downsample is not None:
            identity = self.downsample(identity)

        out = out + identity
        out = F.relu(out, inplace=True)
        return out


# -----------------------------
# 2) ResNet-18 backbone (1D): 4 stages, each stage has 2 blocks
# -----------------------------
class ResNet18_1D(nn.Module):
    

    def __init__(
        self,
        in_channels: int = 1,
        num_classes: int = 2,
        hidden_dim: int = 256,
        dropout_p: float = 0.3,
        use_stats: bool = False, # whether to use the 5 stats features
        stats_dim: int = 5, # number of stats features (if used)
    ):
        super().__init__()
        self.use_stats = use_stats
        self.stats_dim = stats_dim

       
        
        self.stem_conv = nn.Conv1d(in_channels, 64, kernel_size=7, stride=2, padding=3, bias=False)
        self.stem_bn = nn.BatchNorm1d(64)
        self.stem_pool = nn.MaxPool1d(kernel_size=3, stride=2, padding=1)

       
       
        self.in_planes = 64
        # the 4 stages
        self.layer1 = self._make_layer(out_channels=64,  blocks=2, stride=1)
        self.layer2 = self._make_layer(out_channels=128, blocks=2, stride=2)
        self.layer3 = self._make_layer(out_channels=256, blocks=2, stride=2)
        self.layer4 = self._make_layer(out_channels=512, blocks=2, stride=2)

        # Pooling -> embedding
       
        self.avgpool = nn.AdaptiveAvgPool1d(1)

        # FC head (2 FC layers)
       
        fc1_in = embedding_dim + (stats_dim if use_stats else 0)

        self.fc1 = nn.Linear(fc1_in, hidden_dim)
        self.drop = nn.Dropout(p=dropout_p)
        self.fc2 = nn.Linear(hidden_dim, num_classes)

        # init like typical ResNet
        self._init_weights()

    def _make_layer(self, out_channels: int, blocks: int, stride: int):
        layers = []
        # First block may downsample via stride
        layers.append(BasicBlock1D(self.in_planes, out_channels, stride=stride))
        self.in_planes = out_channels * BasicBlock1D.expansion
        # Remaining blocks keep stride=1
        for _ in range(1, blocks):
            layers.append(BasicBlock1D(self.in_planes, out_channels, stride=1))
        return nn.Sequential(*layers)

    def _init_weights(self):
        
        for m in self.modules():
            if isinstance(m, nn.Conv1d):
                nn.init.kaiming_normal_(m.weight, mode="fan_out", nonlinearity="relu")
            elif isinstance(m, nn.BatchNorm1d):
                nn.init.ones_(m.weight)
                nn.init.zeros_(m.bias)

    def forward(self, x, stats=None):
       
        # Stem
        x = self.stem_conv(x)
        x = self.stem_bn(x)
        x = F.relu(x, inplace=True)
        x = self.stem_pool(x)

        # Stages
        x = self.layer1(x)
        x = self.layer2(x)
        x = self.layer3(x)
        x = self.layer4(x)

        # Adaptive pooling -> (B, 512, 1) -> flatten -> (B, 512)
        x = self.avgpool(x)
        x = torch.flatten(x, 1)

        # Optionally concatenate stats features
        if self.use_stats:
            if stats is None:
                raise ValueError("use_stats=True but stats=None was passed to forward().")
            if stats.dim() != 2 or stats.size(1) != self.stats_dim:
                raise ValueError(f"stats must be (B,{self.stats_dim}), got {tuple(stats.shape)}.")
            x = torch.cat([x, stats], dim=1)

        # FC head: FC1 + ReLU + Dropout, then FC2 -> logits
        x = self.fc1(x)
        x = F.relu(x, inplace=True)
        x = self.drop(x)
        logits = self.fc2(x)
        return logits




# -----------------------------
# Metrics (same as your CNN file)
# -----------------------------
def precision_recall_f1(y_true, y_pred):
    y_true = np.asarray(y_true)
    y_pred = np.asarray(y_pred)

    tp = np.sum((y_true == 1) & (y_pred == 1))
    fp = np.sum((y_true == 0) & (y_pred == 1))
    fn = np.sum((y_true == 1) & (y_pred == 0))

    precision = tp / (tp + fp) if (tp + fp) else 0.0
    recall    = tp / (tp + fn) if (tp + fn) else 0.0
    f1        = (2 * precision * recall / (precision + recall)) if (precision + recall) else 0.0
    return float(precision), float(recall), float(f1)


# -----------------------------
# 5 stats features (optional)
# Sheet requires first 3 + 2 proposed :contentReference[oaicite:1]{index=1}
# Proposed extras here:
#   - Shannon entropy of bytes (captures "randomness/structure")
#   - Fraction of zero bytes (captures padding / sparsity patterns)
# -----------------------------
def _shannon_entropy_from_counts(counts: np.ndarray) -> float:
    
    # total number of observed bytes in the histogram
    total = counts.sum()

    # handle edge case of no bytes
    if total <= 0:
        return 0.0

    # convert counts into probabilities of each byte value for non-zero bins only
    p = counts[counts > 0].astype(np.float64) / float(total)

    # compute Shannon entropy -p log2(p) summed over all byte values
    # 0 or 1 entropy if all bytes are the same, and 0.5 is complete uncertainty
    return float(-(p * np.log2(p)).sum())


def compute_stats_features(
    X_bytes: np.ndarray,
    pad_value: int = 0,
    padded: bool = True,
) -> np.ndarray:
   
    X = X_bytes
    if X.dtype != np.uint8:
        X = X.astype(np.uint8)

    N, M = X.shape
    feats = np.zeros((N, 5), dtype=np.float32)

    for i in range(N):
        row = X[i]

        # "total bytes before trimming/padding" requirement :contentReference[oaicite:2]{index=2}
        # If you only have padded arrays, best practical approximation:
        # count non-pad bytes when padding exists, else M.
        if padded:
            total_bytes = int(np.sum(row != pad_value))
            if total_bytes == 0:
                total_bytes = M
        else:
            total_bytes = M
        counts = np.bincount(row, minlength=256)

        # top-2 frequencies
        top2 = np.sort(counts)[-2:]
        top1_count = int(top2[-1])
        top2_count = int(top2[-2])

        ent = _shannon_entropy_from_counts(counts)

        # this tells the model how much of the fixed-length input is artificial. 
        frac_pad = float(counts[pad_value]) / float(M) if M > 0 else 0.0

        # store features
        feats[i, 0] = float(total_bytes)
        feats[i, 1] = float(top1_count)
        feats[i, 2] = float(top2_count)
        feats[i, 3] = float(ent)
        feats[i, 4] = float(frac_pad)

    return feats


# -----------------------------
# Data split for ResNet (PyTorch)
# -----------------------------
def split_training_and_test_resnet(ds, labels, train_indices_fold, test_indices_fold, M: int):
   
    X_train = ds[train_indices_fold]
    X_test  = ds[test_indices_fold]

    # ensure length M (truncate or pad)
    def _fix_len(X):
        if X.shape[1] == M:
            return X
        if X.shape[1] > M:
            return X[:, :M]
        # pad with 0
        pad = np.zeros((X.shape[0], M - X.shape[1]), dtype=X.dtype)
        return np.concatenate([X, pad], axis=1)

    X_train = _fix_len(X_train)
    X_test  = _fix_len(X_test)

    # keep raw bytes for stats (uint8)
    X_train_raw = X_train.astype(np.uint8, copy=False)
    X_test_raw  = X_test.astype(np.uint8, copy=False)

    # normalize to [0,1] and make (N,1,M) to be tensor
    X_train_t = (X_train.astype(np.float32) / 255.0)[:, None, :]
    X_test_t  = (X_test.astype(np.float32) / 255.0)[:, None, :]

    y_train = labels[train_indices_fold].astype(np.int64)
    y_test  = labels[test_indices_fold].astype(np.int64)

    return X_train_t, X_test_t, y_train, y_test, X_train_raw, X_test_raw


# -----------------------------
# Train/Eval ResNet
# -----------------------------
def train_resnet(
    model: torch.nn.Module,
    train_loader: DataLoader,
    val_loader: DataLoader,
    epochs: int = 15,
    lr: float = 1e-3,
    weight_decay: float = 0.0,
    patience: int = 3,
    device: str = "cuda",
):
    # using Adam optimizer
    opt = torch.optim.Adam(model.parameters(), lr=lr, weight_decay=weight_decay)

    # best values tracker
    best_val = float("inf")
    best_state = None

    # Counts how many consecutive epochs did not improve validation loss.
    bad = 0

    # epochs excution
    for _epoch in range(epochs):

        # training
        model.train()
        for batch in train_loader:
            # unpack batch depending on whether stats features exist
            if len(batch) == 2:
                xb, yb = batch
                sb = None # statistics features batch
            else:
                xb, sb, yb = batch

            # move to device
            xb = xb.to(device, non_blocking=True)
            yb = yb.to(device, non_blocking=True)
            if sb is not None:
                sb = sb.to(device, non_blocking=True)
            # forward + backward + optimize steps
            opt.zero_grad(set_to_none=True)
            logits = model(xb, stats=sb) if sb is not None else model(xb)
            loss = F.cross_entropy(logits, yb)
            loss.backward()
            opt.step()

        # validation
        model.eval()
        val_loss = 0.0
        n = 0
        with torch.no_grad():
            for batch in val_loader:
                # same unpacking logic
                if len(batch) == 2:
                    xb, yb = batch
                    sb = None
                else:
                    xb, sb, yb = batch
                xb = xb.to(device, non_blocking=True)
                yb = yb.to(device, non_blocking=True)
                if sb is not None:
                    sb = sb.to(device, non_blocking=True)

                # forward pass + loss computation
                logits = model(xb, stats=sb) if sb is not None else model(xb)
                loss = F.cross_entropy(logits, yb)
                bs = xb.size(0)
                val_loss += float(loss) * bs
                n += bs

        # average validation loss, and avoid division by zero
        val_loss = val_loss / max(n, 1)

        # early stopping check 
        if val_loss < best_val:
            best_val = val_loss
            best_state = {k: v.detach().cpu().clone() for k, v in model.state_dict().items()}
            bad = 0
        else:
            bad += 1
            if bad >= patience:
                break
    # restore best weights at the end from checkpoints. 
    if best_state is not None:
        model.load_state_dict(best_state)

# -----------------------------

def evaluate_resnet(
    model: torch.nn.Module,
    test_loader: DataLoader,
    device: str = "cuda",
):
    model.eval()
    y_true = []
    y_pred = []

    with torch.no_grad():
        for batch in test_loader:
            # same unpacking logic
            if len(batch) == 2:
                xb, yb = batch
                sb = None
            else:
                xb, sb, yb = batch

            xb = xb.to(device, non_blocking=True)
            if sb is not None:
                sb = sb.to(device, non_blocking=True)

            logits = model(xb, stats=sb) if sb is not None else model(xb)
            preds = torch.argmax(logits, dim=1).cpu().numpy()

            y_pred.append(preds)
            y_true.append(yb.numpy())
    # concatenate all batches
    y_true = np.concatenate(y_true)
    y_pred = np.concatenate(y_pred)

    # compute metrics
    p, r, f1 = precision_recall_f1(y_true, y_pred)
    return p, r, f1, y_pred


# -----------------------------
# One fold execution
# -----------------------------
def execute_fold_resnet(
    fold_idx: int,
    binary_numeric_labels: np.ndarray,
    scenario: int,
    param: int,
    train_idx,
    test_idx,
    M: int,
    use_stats: bool = False,
    stats_dim: int = 5,
    epochs: int = 15,
    batch_size: int = 256,
    lr: float = 1e-3,
    dropout_p: float = 0.3,
    hidden_dim: int = 256,
    weight_decay: float = 0.0,
    device: str | None = None,
    debug_subset: bool = True,  # mimic your CNN "first+last 1000"
):
    if device is None:
        device = "cuda" if torch.cuda.is_available() else "cpu"

    # ---- load dataset ----
    if param != 0:
        ds = np.load(f"datasets/re_bytes_{param}.npy")
        print(f"Experiment: ResNet classifier, fold {fold_idx}, scenario {scenario}, RE{param}, stats={use_stats}")
    else:
        ds = np.load("datasets/raw_bytes.npy")
        print(f"Experiment: ResNet classifier, fold {fold_idx}, scenario {scenario}, RAW, stats={use_stats}")

    # split -> normalize -> (N,1,M)
    X_train, X_test, y_train, y_test, X_train_raw, X_test_raw = split_training_and_test_resnet(
        ds, binary_numeric_labels, np.array(train_idx), np.array(test_idx), M=M
    )

    # optional small subset (same spirit as your CNN demo)
    if debug_subset and len(X_train) > 2000:
        first = np.arange(0, 1000)
        last  = np.arange(len(X_train) - 1000, len(X_train))
        sel = np.concatenate([first, last])
        X_train = X_train[sel]
        y_train = y_train[sel]
        X_train_raw = X_train_raw[sel]

    # stats (if enabled)
    if use_stats:
        S_train = compute_stats_features(X_train_raw)  # (Ntr,5)
        S_test  = compute_stats_features(X_test_raw)   # (Nte,5)
        S_train_t = torch.from_numpy(S_train)
        S_test_t  = torch.from_numpy(S_test)
    else:
        S_train_t = None
        S_test_t  = None

    # tensors
    X_train_t = torch.from_numpy(X_train)
    y_train_t = torch.from_numpy(y_train)
    X_test_t  = torch.from_numpy(X_test)
    y_test_t  = torch.from_numpy(y_test)

    # train/val split
    n = len(X_train_t)
    perm = torch.randperm(n)
    val_n = max(1, int(0.1 * n))
    val_idx = perm[:val_n]
    tr_idx  = perm[val_n:]

    if use_stats:
        train_ds = TensorDataset(X_train_t[tr_idx], S_train_t[tr_idx], y_train_t[tr_idx])
        val_ds   = TensorDataset(X_train_t[val_idx], S_train_t[val_idx], y_train_t[val_idx])
        test_ds  = TensorDataset(X_test_t, S_test_t, y_test_t)
    else:
        train_ds = TensorDataset(X_train_t[tr_idx], y_train_t[tr_idx])
        val_ds   = TensorDataset(X_train_t[val_idx], y_train_t[val_idx])
        test_ds  = TensorDataset(X_test_t, y_test_t)

    train_loader = DataLoader(train_ds, batch_size=batch_size, shuffle=True, drop_last=False)
    val_loader   = DataLoader(val_ds, batch_size=batch_size, shuffle=False, drop_last=False)
    test_loader  = DataLoader(test_ds, batch_size=batch_size, shuffle=False, drop_last=False)

    # model
    model = ResNet18_1D(
        in_channels=1,
        num_classes=2,
        hidden_dim=hidden_dim,
        dropout_p=dropout_p,
        use_stats=use_stats,
        stats_dim=stats_dim,
    ).to(device)

    train_resnet(
        model=model,
        train_loader=train_loader,
        val_loader=val_loader,
        epochs=epochs,
        lr=lr,
        weight_decay=weight_decay,
        patience=3,
        device=device,
    )

    # eval
    p, r, f1, _ = evaluate_resnet(model=model, test_loader=test_loader, device=device)
    return p, r, f1


# -----------------------------
# Run ResNet for a scenario + representation
# -----------------------------
def run_resnet_for_scenario(
    scenario: int,
    prefix_for_files: str,   # "raw" or "re"
    global_label_encoder,
    keep_indices,
    param: int,
    M: int,
    use_stats: bool = False,
    stats_dim: int = 5,
    epochs: int = 15,
    batch_size: int = 256,
    lr: float = 1e-3,
    dropout_p: float = 0.3,
    hidden_dim: int = 256,
    weight_decay: float = 0.0,
):
    labels = np.load(f"datasets/{prefix_for_files}_labels.npy")

    train_indices, test_indices = load_k_fold_results(
        f"k_fold_results/k_fold_s{scenario}_{prefix_for_files}.json"
    )

    if param != 0:
        train_indices, test_indices = deduplicate_folds(train_indices, test_indices, keep_indices)

    numeric_labels = encode_labels(global_label_encoder, labels)
    binary_numeric_labels = np.where(numeric_labels == 0, 0, 1)

    k = len(train_indices)

    # Fine-tuning hyperparameters
    epochss = [15, 20, 25,30]
    batch_sizes = [256, 512, 1024]
    lrs = [1e-3, 5e-4, 1e-4]
    dropout_ps = [0.3, 0.5, 0.7]
    
    best_epoch = -1
    best_bs = -1
    best_lr = -1
    best_dps = -1 

    best_f1 = 0
    best_percision = 0
    best_recall = 0

    for e in epochss:
        for bs in batch_sizes:
            for learning_rate in lrs:
                for dp in dropout_ps:          
                    print(f'running with current parameters: epochs={e}, batch_size={bs}, learning_rate={learning_rate}, dropout_p={dp}')        
                    precisions, recalls, f1s = [], [], []
                    for fold_idx in range(k):
                        if len(train_indices[fold_idx]) == 0 or len(test_indices[fold_idx]) == 0:
                            continue

                        p, r, f1 = execute_fold_resnet(
                            fold_idx=fold_idx,
                            binary_numeric_labels=binary_numeric_labels,
                            scenario=scenario,
                            param=param,
                            train_idx=train_indices[fold_idx],
                            test_idx=test_indices[fold_idx],
                            M=M,
                            use_stats=use_stats,
                            stats_dim=stats_dim,
                            epochs=e,
                            batch_size=bs,
                            lr=learning_rate,
                            dropout_p=dp,
                            hidden_dim=hidden_dim,
                            weight_decay=weight_decay,
                        )

                        precisions.append(p)
                        recalls.append(r)
                        f1s.append(f1)
                        print(f"[Fold {fold_idx}] P={p:.4f} R={r:.4f} F1={f1:.4f}")
                    avg_f1 = float(np.mean(f1s))
                    if avg_f1 > best_f1:
                        best_f1 = avg_f1
                        best_percision = float(np.mean(precisions))
                        best_recall = float(np.mean(recalls))
                        best_epoch = e
                        best_bs = bs
                        best_lr = learning_rate
                        best_dps = dp
    # if len(precisions) == 0:
    #     return 0.0, 0.0, 0.0
    print(f'best parameters are: {best_epoch}, {best_bs}, {best_lr}, {best_dps}')
    print(f'best results: P={best_percision:.4f} R={best_recall:.4f} F1={best_f1:.4f}')

    return best_percision, best_recall, best_f1  
    return float(np.mean(precisions)), float(np.mean(recalls)), float(np.mean(f1s))


def run_resnet_classifier_on_dataset(
    scenario: int,
    prefix: str,               # "raw", "re5", "re10", "re15" (only used for naming)
    param: int,                # 0, 5, 10, 15
    global_label_encoder,
    M: int,
    use_stats: bool = False,
    out_csv: str = "results/resnet_summary.csv",
):
    if param != 0:
        keep_indices = get_keep_indices_from_fold0_ae(f"datasets/re_bytes_{param}.npy")
        prefix_for_files = "re"
    else:
        keep_indices = []
        prefix_for_files = "raw"

    avg_p, avg_r, avg_f1 = run_resnet_for_scenario(
        scenario=scenario,
        prefix_for_files=prefix_for_files,
        global_label_encoder=global_label_encoder,
        keep_indices=keep_indices,
        param=param,
        M=M,
        use_stats=use_stats,
        stats_dim=5,
    )

    out_dir = Path("results")
    out_dir.mkdir(exist_ok=True)

    summary_path = Path(out_csv)
    write_header = not summary_path.exists()

    representation = "raw" if param == 0 else f"re{param}"

    with summary_path.open("a", newline="") as f:
        w = csv.writer(f)
        if write_header:
            w.writerow(["scenario", "representation", "M", "use_stats", "avg_precision", "avg_recall", "avg_f1"])
        w.writerow([scenario, representation, M, int(use_stats), f"{avg_p:.6f}", f"{avg_r:.6f}", f"{avg_f1:.6f}"])

    return avg_p, avg_r, avg_f1


# -----------------------------
# The function you asked for:
# similar to run_experiment_cnn_classifier
# -----------------------------
def run_experiment_resnet_classifier(global_label_encoder, M_raw, M_re):
   
    prefixes = ["raw", "re5", "re10", "re15"]
    params = [0, 5, 10, 15]

    results = {}
    for scenario in (2, 3):
        results[scenario] = {}
        for prefix, param in zip(prefixes, params):
            M = M_raw if param == 0 else M_re

            # no-stats
            key0 = f"{prefix}_nostats"
            results[scenario][key0] = run_resnet_classifier_on_dataset(
                scenario=scenario,
                prefix=prefix,
                param=param,
                global_label_encoder=global_label_encoder,
                M=M,
                use_stats=False,
                out_csv="results/resnet_summary.csv",
            )

            # with-stats
            key1 = f"{prefix}_withstats"
            results[scenario][key1] = run_resnet_classifier_on_dataset(
                scenario=scenario,
                prefix=prefix,
                param=param,
                global_label_encoder=global_label_encoder,
                M=M,
                use_stats=True,
                out_csv="results/resnet_summary.csv",
            )

    return results

\end{lstlisting}


\section{Implementation of the GAN-based Anomaly Detector}

This appendix presents the implementation of the Generative Adversarial Network (GAN) used for anomaly detection in industrial control system (ICS) network traffic. The full implementation is provided in Listing~\ref{lst:gan_code}. The code is implemented in Python using the PyTorch deep learning framework.

The implemented GAN consists of two main components: a \textit{generator} and a \textit{discriminator}. The generator learns to produce synthetic packet representations, while the discriminator attempts to distinguish between real network packets and generated samples. Through this adversarial training process, the discriminator learns meaningful representations of normal network traffic, which can later be used for anomaly detection.

\subsection{Discriminator Architecture}

The discriminator is implemented as a modular model composed of two parts: a feature extractor and a classifier. This design allows the model to expose intermediate feature representations that are later used for feature matching during GAN training.

\begin{itemize}

\item \textbf{Feature Extractor.}  
The feature extractor consists of nine convolutional layers using two-dimensional convolutions. The input data is represented as a tensor of shape $(batch\_size, 1, m, n)$, where $m$ denotes the number of packets in a sequence and $n$ denotes the number of bytes per packet. 

Each convolutional layer is followed by a ReLU activation function. Dropout regularization is applied after most convolutional layers in order to improve generalization and reduce overfitting. The number of filters gradually increases across layers, enabling the network to learn hierarchical representations of packet structures. After the convolutional layers, an adaptive average pooling layer reduces the spatial dimensions and produces a compact feature representation.

\item \textbf{Classifier.}  
The extracted feature vector is passed to a fully connected classifier composed of four dense layers with ReLU activation functions. The final layer outputs a single value representing the discriminator score, which indicates how likely the input sample is to belong to the real data distribution.

\item \textbf{Wrapper Model.}  
The final discriminator model combines the feature extractor and the classifier into a single module. This wrapper enables two types of forward passes:
\begin{itemize}
\item extraction of intermediate feature representations used for feature matching loss,
\item computation of the final real/fake score used during adversarial training.
\end{itemize}

\end{itemize}

\subsection{Generator Architecture}

The generator network is designed to transform random noise vectors into synthetic packet representations that resemble real ICS network traffic.

The generator receives a noise vector $z \sim \mathcal{N}(0,1)$ with dimensionality $m \times n$. The architecture consists of two main stages:

\begin{itemize}

\item \textbf{Fully Connected Layers.}  
The first stage contains four fully connected layers with LeakyReLU activation functions. These layers expand the noise vector into a high-dimensional representation suitable for convolutional processing.

\item \textbf{Convolutional Layers.}  
The second stage consists of eight convolutional layers that progressively transform the intermediate representation into a tensor with the same spatial dimensions as the input packet representation $(1, m, n)$. The first four convolutional layers are preceded by upsampling operations, enabling the generator to construct spatial structure in the generated packets.

The final convolutional layer uses a sigmoid activation function to ensure that the generated values lie within the normalized range of the packet data.

\end{itemize}

\subsection{Model Integration}

During training, the generator and discriminator are optimized in an adversarial manner. The discriminator learns to differentiate between real network packets and generated samples, while the generator attempts to produce samples that the discriminator cannot distinguish from real data. In addition to the standard adversarial objective, the discriminator features are used for feature matching in order to stabilize training.

The complete implementation of the GAN components used in our experiments is shown below.

\begin{lstlisting}[language=Python, caption={Implementation of the GAN discriminator and generator models}, label={lst:gan_code}]

import torch
import torch.nn as nn


class Generator(nn.Module):
    def __init__(self, m, n, kernel_size=3):
        super().__init__()

        self.m = m
        self.n = n
        self.z_dim = m * n

        # -------- Stage 1: Fully Connected layers --------
        self.fc = nn.Sequential(
            nn.Linear(self.z_dim, 512),
            nn.LeakyReLU(0.2),

            nn.Linear(512, 1024),
            nn.LeakyReLU(0.2),

            nn.Linear(1024, 2048),
            nn.LeakyReLU(0.2),

            nn.Linear(2048, 256 * m * n),
            nn.LeakyReLU(0.2)
        )

        # -------- Stage 2 & 3: Convolutional layers --------
        self.conv_layers = nn.Sequential(
            # Upsampling + Conv (1)
            nn.Upsample(scale_factor=1),  # placeholder (keeps shape explicit)
            nn.Conv2d(256, 256, kernel_size=kernel_size, padding=1),
            nn.LeakyReLU(0.2),

            # Upsampling + Conv (2)
            nn.Upsample(scale_factor=1),
            nn.Conv2d(256, 256, kernel_size=kernel_size, padding=1),
            nn.LeakyReLU(0.2),

            # Upsampling + Conv (3)
            nn.Upsample(scale_factor=1),
            nn.Conv2d(256, 128, kernel_size=kernel_size, padding=1),
            nn.LeakyReLU(0.2),

            # Upsampling + Conv (4)
            nn.Upsample(scale_factor=1),
            nn.Conv2d(128, 64, kernel_size=kernel_size, padding=1),
            nn.LeakyReLU(0.2),

            # Conv (5)
            nn.Conv2d(64, 64, kernel_size=kernel_size, padding=1),
            nn.LeakyReLU(0.2),

            # Conv (6)
            nn.Conv2d(64, 32, kernel_size=kernel_size, padding=1),
            nn.LeakyReLU(0.2),

            # Conv (7)
            nn.Conv2d(32, 16, kernel_size=kernel_size, padding=1),
            nn.LeakyReLU(0.2),

            # Conv (8) - output layer
            nn.Conv2d(16, 1, kernel_size=kernel_size, padding=1),
            nn.Sigmoid()
        )

    def forward(self, z):
        
        x = self.fc(z)
        x = x.view(z.size(0), 256, self.m, self.n)
        x = self.conv_layers(x)
        return x


import torch
import torch.nn as nn


class DiscriminatorClassifier(nn.Module):
    def __init__(self, feature_dim):
       
        super().__init__()

        # typical architecture for binary classification
        self.fc_layers = nn.Sequential(
            nn.Linear(feature_dim, 512), # map the input feature vector to 512-dim
            nn.ReLU(), # add non-linearity

            nn.Linear(512, 256), # map to 256-dim
            nn.ReLU(), # add non-linearity for more complex decision boundary

            nn.Linear(256, 128),
            nn.ReLU(),

            nn.Linear(128, 1),
            # nn.Sigmoid() # output probability of being real (1) or fake (0)
        )

    def forward(self, features):
        return self.fc_layers(features)

import torch
import torch.nn as nn
from  discriminator_classifier import DiscriminatorClassifier
from discriminator_feature_extractor import DiscriminatorFeatures



class Discriminator(nn.Module):
    def __init__(self, input_shape):
        super().__init__()

        self.feature_extractor = DiscriminatorFeatures()

        # Compute feature dimension dynamically
        with torch.no_grad(): # disable gradient calculation because this step is not part of training

            # create fake input tensor filled with zeros
            dummy = torch.zeros(1, *input_shape) # 1 is batch size, *input_shape unpacks the input shape tuple

            # runs this dummy input through the feature extractor to get the feature dimension
            feature_dim = self.feature_extractor(dummy).shape[1] # number of features per sample. 

        # define a classifier using the computed feature dimension
        self.classifier = DiscriminatorClassifier(feature_dim)

    def extract_features(self, x):
        
        return self.feature_extractor(x)

    def forward(self, x):
        
        features = self.extract_features(x)
        score = self.classifier(features)
        return score
\end{lstlisting}

\section{Implementation of the N-gram Based Detection Method}

This appendix presents the implementation of the N-gram based anomaly detection method used in our experiments. The complete source code is shown in Listing~\ref{lst:ngram_code}. The implementation is written in Python and follows the design requirements described in the project specification.

The core idea of the method is to model the byte-level structure of network packets using N-grams. During the training phase, the model observes all N-grams appearing in normal packets and records their frequencies. During testing, packets are scored based on how many of their N-grams were previously observed during training.

\subsection{Bloom Filter for Efficient Storage}

To efficiently store the set of observed N-grams, the implementation uses a Bloom filter. A Bloom filter is a probabilistic data structure that allows fast membership queries while using very little memory. 

The size of the Bloom filter and the number of hash functions are computed automatically using the standard formulas:

\[
m = -\frac{n \ln p}{(\ln 2)^2}, \qquad
k = \frac{m}{n} \ln 2
\]

where $n$ is the expected number of stored N-grams and $p$ is the desired false positive rate. The implementation uses double hashing based on MurmurHash3 and BLAKE2b in order to generate multiple hash values efficiently.

\subsection{Training Phase}

During the training phase, the algorithm performs the following steps:

\begin{enumerate}

\item Extract all N-grams from each packet in the training dataset.

\item Count the frequency of each N-gram across the training data.

\item Insert each distinct N-gram into the Bloom filter.

\item Compute normalized weights for each N-gram using:

\[
w_i = \frac{\text{frequency of } n\text{-gram}_i}{\text{total number of observed } n\text{-grams}}
\]

\end{enumerate}

These weights represent how frequently each N-gram appears in the training data and are later used during packet scoring.

\subsection{Packet Scoring}

For each packet in the testing phase, the algorithm extracts its N-grams and checks whether they were observed during training using the Bloom filter.

The score of a packet is computed as:

\[
Score = \frac{\sum_{i=1}^{T} w_i \cdot N_{\text{seen}_i}}{T}
\]

where:

\begin{itemize}
\item $T$ is the total number of N-grams in the packet,
\item $w_i$ is the normalized weight of the $i$-th N-gram,
\item $N_{\text{seen}_i}$ indicates whether the N-gram was observed during training.
\end{itemize}

Packets with scores below a threshold are classified as attacks, while packets with higher scores are considered normal.

\subsection{Threshold Selection}

To determine the optimal decision threshold, the implementation evaluates different threshold values on a validation set. The threshold is selected by maximizing the F1-score computed from precision and recall values.

\subsection{Evaluation}

The final model is evaluated on a separate testing dataset. The implementation computes standard classification metrics, including accuracy, precision, recall, and F1-score.

The full implementation of the N-gram detection method is provided below.

\begin{lstlisting}[language=Python, caption={Implementation of the N-gram based anomaly detection method}, label={lst:ngram_code}]

TESTING = False
import sys
import numpy as np
import math
import os
import pickle
import hashlib
import mmh3  # MurmurHash3 library for fast hashing

# ----------------------------------------------------------
# BLOOM FILTER
# ----------------------------------------------------------
class BloomFilter:
   
    def __init__(self, false_positive_rate=0.01, number_of_ngrams=1000):
        

        if number_of_ngrams <= 0 or false_positive_rate <=0 or false_positive_rate >=1:
            raise ValueError("Number of n-grams must be positive and false_positive_rate must be in (0,1) range")
        
        ln_p = math.log(false_positive_rate) # math.log is natural logarithm (ln)
        self.required_number_of_bits = int(-(number_of_ngrams * ln_p) / (math.log(2) ** 2)) # optimal size in bits The Bloom filter size formula is \(m=-(n*\ln (p))/(\ln (2)^{2})\)
        self.number_of_hash_functions = math.ceil((self.required_number_of_bits / number_of_ngrams) * math.log(2)) # k = (m / n) * ln(2)
        self.byte_array = bytearray(self.required_number_of_bits // 8 + 1)  # +1 to handle any remainder bits
        # bytes array is an efficient way to store bits in python, each byte has 8 bits, so number of bytes = ceil(number of bits / 8)
        # I could have used also boolean array, but 1 boolean = 1 byte, so it is less efficient in terms of space.
        # number of hash functions affect the probability of false positives, less hashes = more false positives, more hashes = less false positives but slower performance.


    def _set_certain_bit(self, bit_index):
       

        byte_index = bit_index // 8
        bit_position = bit_index % 8
        self.byte_array[byte_index] |= (1 << bit_position) # set the bit using orwise operation by shifting 1 to the left by bit_position
    
    def _get_certain_bit(self, bit_index):
       
        byte_index = bit_index // 8
        bit_position = bit_index % 8
        return (self.byte_array[byte_index] >> bit_position) & 1 # get the bit using andwise operation by shifting right by bit_position till the 0th position which include my 1.

        
    def _hashes(self, data: bytes):
        
        # First hash: MurmurHash3 (extremely fast, 64-bit output)
        h1, _ = mmh3.hash64(data, seed=0, signed=False)

        # Second hash: BLAKE2b (strong entropy, 64-bit output)
        blake = hashlib.blake2b(data, digest_size=8).digest()
        h2 = int.from_bytes(blake, "big") | 1  # ensure h2 is odd, to avoid cycles, and try to have all k hash values different

        # Produce k hash values using double hashing
        for i in range(self.number_of_hash_functions):
            yield (h1 + i * h2) % self.required_number_of_bits # return one value now, but remember where we left off to continue later.


    def add(self, ngram):
       
        # I wanted to check if the element already exist, but I thought that due to collisions, it is better to just set the bits again.
        for h in self._hashes(ngram):
            self._set_certain_bit(h)

    def __contains__(self, ngram):
       
        return all(self._get_certain_bit(h) for h in self._hashes(ngram)) # all check that all bits are set to 1, which means the n-gram is probably in the filter.
        # I said probably because of false positives, and collisions. 

# ----------------------------------------------------------
# N-GRAM EXTRACTION
# ----------------------------------------------------------
def extract_ngrams(byte_seq, n):
    return [byte_seq[i:i+n] for i in range(len(byte_seq) - n + 1)]


# ----------------------------------------------------------
# TRAINING PHASE
# ----------------------------------------------------------
def train_ngram_models(train_packets, n):
    

    ngram_count = {}
    total_seen = 0 # total number of distinct n-grams seen during training

    print(f'starting to train n-gram model...\n number of training packets: {len(train_packets)}')
    for pkt in train_packets:
        # for i in range(2,n+1):# generating many n-grams till the maximum n (include n).
            # pass only first element in the tuple (packet, label)
        grams = extract_ngrams(pkt[0], n)
        for g in grams:
            g = bytes(g)
            ngram_count[g] = ngram_count.get(g, 0) + 1
            total_seen += 1
    
    # add n-grams to bloom filter
    print('starting to add n-grams to bloom filter...')
    bloom = BloomFilter(false_positive_rate=0.01, number_of_ngrams=len(train_packets)*1000) # assuming each packet has 1000 distinct n-grams on average
    for curr_gram in ngram_count.keys():
        bloom.add(curr_gram)

    # normalized weights
    print('starting to normalize weights...')
    weights = {curr_gram: c / total_seen for curr_gram, c in ngram_count.items()} # I count the weights only to certan n-grams seen in training data
    # i.e. 0100 was seen 5 times, and total 2-grams are 100, then weight of 0100 = 5/100 = 0.05
    # then I will do 010011 for example which is 3-gram, so I will now count total seen different from the 100 of 2-grams, I will check how many 3-grams only.
            
    # sample from weights
    print(f'samples of weights are: {list(weights.items())[:20]}')

    return bloom, weights, ngram_count


# ----------------------------------------------------------
# SCORE A PACKET
# ----------------------------------------------------------
def score_packet(packet, bloom, weights, n): 
    
    T = 0 # represent total number of n-grams in the current packet.
    score_sum = 0
    # for i in range(2,n+1):# generating many n-grams till the maximum n (include n).
        # # print(f'current packet is: {packet}')
    grams = extract_ngrams(packet, n)
    T += len(grams) # total number of n-grams in the current packet.
    if len(grams) == 0:
        return 0

    
    for g in grams: # iterate for each gram in the current packet
        g = bytes(g)
        if g in bloom: # if it was seen during training
            w = weights.get(g, 0) # get its weight from the training phase
            # N_seen_i = ngram_count.get(g, 0) # get its count from the training phase
            score_sum += w  # compute score contribution: (N_seen_i * wi)
            # when I include N_seen_i, the score becomes > 1, and I think this is because we double count the counts.

    final_score = score_sum / T if T > 0 else 0
    return final_score


# ----------------------------------------------------------
# SEARCH FOR BEST THRESHOLD
# ----------------------------------------------------------

def compute_metrics(y_true, y_pred): # recheck this logic.
   
    tp = sum(1 for yt, yp in zip(y_true, y_pred) if yt == 'attack' and yp == 'attack')
    tn = sum(1 for yt, yp in zip(y_true, y_pred) if yt == 'normal' and yp == 'normal')
    fp = sum(1 for yt, yp in zip(y_true, y_pred) if yt == 'normal' and yp == 'attack')
    fn = sum(1 for yt, yp in zip(y_true, y_pred) if yt == 'attack' and yp == 'normal')
    accuracy = (tp + tn) / (tp + tn + fp + fn) if (tp + tn + fp + fn) > 0 else 0
    precision = tp / (tp + fp) if (tp + fp) > 0 else 0
    recall = tp / (tp + fn) if (tp + fn) > 0 else 0
    f1_score = 2 * (precision * recall) / (precision + recall) if (precision + recall) > 0 else 0

    return accuracy, precision, recall, f1_score


def find_best_threshold(test_set, scores): # recheck this logic.
   

    best_acc = -1
    best_prec = -1
    best_rec = -1
    best_f1 = -1
    best_t = None

    # -------------------------------------------------
    # 1. Compute scores for all test packets ONCE only
    # -------------------------------------------------
    # scores = [score_packet(pkt, bloom, weights, n) for pkt, _ in test_set]
    y_true = [label for _, label in test_set]

    # -------------------------------------------------
    # 2. Sweep thresholds between 0 and 1
    # -------------------------------------------------
    # # print how many normal, and how many attack labels are in y_true
    num_normal = sum(1 for label in y_true if label == 'normal')
    num_attack = sum(1 for label in y_true if label == 'attack')
    print(f'Number of normal packets: {num_normal}, Number of attack packets: {num_attack}\n total: {len(y_true)}')


    highest_prec = -1
    highest_rec = -1
    highest_acc = -1
    highest_f1 = -1
    for t in np.linspace(0, 1, 1000):
        # print(f'Evaluating threshold: {t:.4f}')
        # 3. Predict using threshold t
        y_pred = ['normal' if s >= t else 'attack' for s in scores] 

        # 4. Compute metrics
        acc, prec, rec, f1_score = compute_metrics(y_true, y_pred)

        # 5. Compare with previous best
        if prec > highest_prec:
            highest_prec = prec
        if rec > highest_rec:
            highest_rec = rec
        if acc > highest_acc:
            highest_acc = acc
        if f1_score > highest_f1:
            highest_f1 = f1_score
        
        if f1_score > best_f1: # if I prioritize recall, all of them will be normal, and if I prioritize precision, all of them will be attack.
            best_acc = acc
            best_prec = prec
            best_rec = rec
            best_f1 = f1_score
            best_t = t
    print(f'Highest Precision observed: {highest_prec}'
          f', Highest Recall observed: {highest_rec}'
        f', Highest Accuracy observed: {highest_acc}'
        f', Highest F1 Score observed: {highest_f1}')
    return best_t, best_acc, best_prec, best_rec, best_f1


# ----------------------------------------------------------
# TESTING
# ----------------------------------------------------------
def test_model(test_packets, bloom, weights, n, threshold):
    results = []
    for i, pkt in enumerate(test_packets):
        s = score_packet(pkt[0], bloom, weights, n)
        label = "normal" if s >= threshold else "attack"
        results.append((i, s, label))
    # compute metrics for the test set
    y_true = [label for _, label in test_packets]
    y_pred = [label for _, _, label in results]
    acc, prec, rec, f1_score = compute_metrics(y_true, y_pred)
    print(f"[*] Test Set Results: Accuracy={acc:.4f}, Precision={prec:.4f}, Recall={rec:.4f}, F1 Score={f1_score:.4f}")
    # print how many normal and how many attack packets were detected
    normal_detected = sum(1 for _, _, label in results if label == 'normal')
    attack_detected = sum(1 for _, _, label in results if label == 'attack')
    print(f'Number of NORMAL packets detected: {normal_detected}, Number of ATTACK packets detected: {attack_detected}\n total: {len(results)}')
    print("[*] Sample results (index, score, label):", results[:10])

    return results




def load_all_modules():
    import sys
 

    # print('loading all necessary modules')
    curr_dir = os.path.dirname(os.path.abspath(__file__))
    sheet1_codes_path = os.path.abspath(os.path.join(curr_dir, '..', '..','Assignment1','Abdelaziz_Codes' ,'Sheet1_codes'))
    sys.path.append(sheet1_codes_path)
    

def load_and_label_data():
    normal_pcap_path = "../../DataSets/2017QUT_S7comm/LabelledDataset/20161219132813_control_set"
    attacked_pcap_path  = "../../DataSets/2017QUT_S7comm/LabelledDataset/20161215163606_s7_process_attacks"
    load_all_modules()
    from utilities import generate_bytes_array_from_packet_list

    # check if .npy files already exist, read it if yes, else generate it from pcap
    print(f"[*] Loading normal packets...")
    if os.path.exists("all_packets_control.npy"):
        data_normal_arrays = np.load("all_packets_control.npy", allow_pickle=True)
    else:
        data_normal_arrays = generate_bytes_array_from_packet_list(normal_pcap_path, pad = False, label = 'control')
 

    print(f"[*] Loading attacked packets...")
    data_attacked = []
    if os.path.exists("all_packets_attack.npy"):
        data_attacked = np.load("all_packets_attack.npy", allow_pickle=True)
    else:
        data_attacked = generate_bytes_array_from_packet_list(attacked_pcap_path, pad = False, label = 'attack')
    
    labeled_data_normal =[ ]
    labeled_data_attack =[ ]
    for array in data_normal_arrays:
        for pkt in array:
            labeled_data_normal.append( (pkt, 'normal') ) 

    if TESTING:
        for array in data_attacked[-2]:# this logic should be the same as normal.
            # for pkt in array:
            labeled_data_attack.append( (pkt, 'attack') )
    else:
        for array in data_attacked:
            for pkt in array:
                labeled_data_attack.append( (pkt, 'attack') )
    return labeled_data_normal, labeled_data_attack

def compute_bloom_weights_and_counts(train_packets, n):
    
    print("[*] Training model...")
    if os.path.exists(f'bloom_filter_{n}.pkl') and os.path.exists(f'ngram_weights_{n}.pkl') and os.path.exists(f'ngram_count_training_{n}.pkl'):
        
        with open(f'bloom_filter_{n}.pkl', 'rb') as f:
            bloom = pickle.load(f)
        with open(f'ngram_weights_{n}.pkl', 'rb') as f:
            weights = pickle.load(f)
        with open(f'ngram_count_training_{n}.pkl', 'rb') as f:
            ngram_count_training = pickle.load(f)
        print(f"[*] Loaded pre-trained model from disk. Number of n-grams in model: {len(weights)}")
        print('-------------------------------------')
    else:
        bloom, weights, ngram_count_training = train_ngram_models(train_packets, n)
        # save the bloom filter and weights to disk for future use
        with open(f'bloom_filter_{n}.pkl', 'wb') as f:
            pickle.dump(bloom, f)
        with open(f'ngram_weights_{n}.pkl', 'wb') as f:
            pickle.dump(weights, f)
        # store ngram_count_training for future use
        with open(f'ngram_count_training_{n}.pkl', 'wb') as f:
            pickle.dump(ngram_count_training, f)
        print(f"[*] Training completed. Number of n-grams in model: {len(weights)}")
        print('-------------------------------------')
    return bloom, weights


def split_data(labeled_data_normal, labeled_data_attack, train_ratio=0.6, validation_ratio=0.2):
    
    # split data into training, validaiton and testing sets
    split_idx_normal = int(train_ratio * len(labeled_data_normal)) # [ 60% normal packets for training, 20% for validation, 20% for testing ]
    train_normal, remaining_normal = labeled_data_normal[:split_idx_normal], labeled_data_normal[split_idx_normal:]
    split_idx_normal_val = int(validation_ratio * len(remaining_normal))
    validation_normal, test_normal = remaining_normal[:split_idx_normal_val], remaining_normal[split_idx_normal_val:]

    # split attack data into validation and testing sets only => [ 0% attack packets for training, 80% for validation, 20% for testing ]
    split_idx_attack = int(validation_ratio * len(labeled_data_attack))
    test_attack,validation_attack =  labeled_data_attack[:split_idx_attack], labeled_data_attack[split_idx_attack:]
   
    # assign the sets
    train_packets = train_normal  # only normal packets for training
    validation_packets = validation_normal + validation_attack # both normal and attack packets for validation
    test_packets = test_normal + test_attack # both normal and attack packets for testing
    if TESTING:
        # include 8000 from normal packets and 2000 from attack packets in the train set
        train_packets = train_packets[:80000]
        validation_packets = validation_normal[:20000] + validation_attack[:80000]
        # include 2000 normal packets and 2000 attack packets in the test set
        test_packets = test_normal[:10000] + test_attack[:30000]
        print(f'len(train_packets): {len(train_packets)}')
        print(f'len(validation_packets): {len(validation_packets)}')
        print(f'len(test_packets): {len(test_packets)}')
    return train_packets, test_packets, validation_packets
    
# ----------------------------------------------------------
# MAIN ENTRY (CLI)
# ----------------------------------------------------------
def sheet3_task1():
    # read n from command line arguments
    print('-------------------------------------')
    print("[*] Starting N-gram based Detection Method...")
    
    if len(sys.argv) > 1:
        n = int(sys.argv[1])
        print(f"[*] Using n-gram size: {n}")
    else:
        # print a hint how to call it in command line
        raise ValueError("Please provide n-gram size as a command line argument. Example: python task1_sheet3.py 3")
        # n = 3  # default n-gram size

    # load and label data
    print('[*] Loading and labeling data...')
    labeled_data_normal, labeled_data_attack = load_and_label_data() 
    
    # split data into training and testing sets
    print("[*] Splitting data into training and testing sets...")
    train_packets, test_packets, validation_packets = split_data(labeled_data_normal, labeled_data_attack, train_ratio=0.6, validation_ratio=0.2)


    # compute bloom filter, weights, and ngram counts from training data
    print("[*] Train n-gram model and compute bloom filter, weights, and ngram counts...")
    bloom, weights = compute_bloom_weights_and_counts(train_packets, n)


    print("[*] Scoring validation packets for threshold search...")
    # validation should be on both normal and attack packets
    # normal packets should have high scores, while attack packets should have low scores
    validation_scores = [score_packet(pkt[0], bloom, weights, n)
                    for pkt in validation_packets]
    # save validation scores to a file
    with open(f'validation_scores_{n}.txt', 'w') as f:
        f.write("Index\tScore\tLabel\n")
        for index, (_, label) in enumerate(validation_packets):
            f.write(f"{index}\t{validation_scores[index]:.6f}\t{label}\n")
    print("[*] Sample validation scores:", validation_scores[:10])
    print('-------------------------------------')

    print("[*] Finding optimal threshold...")
    best_t, best_acc, best_prec, best_rec, best_f1 = find_best_threshold(validation_packets,scores=validation_scores)
    # save best threshold to a file
    with open(f'best_metrics_{n}.txt', 'w') as f:
        f.write(f"Best Threshold: {best_t:.6f}\n")
        f.write(f"Accuracy: {best_acc:.6f}\n")
        f.write(f"Precision: {best_prec:.6f}\n")
        f.write(f"Recall: {best_rec:.6f}\n")
        f.write(f"F1 Score: {best_f1:.6f}\n")
    print(f"[+] Optimal threshold on validation found: {best_t:.4f} with Accuracy={best_acc:.4f}, Precision={best_prec:.4f}, Recall={best_rec:.4f}, F1 Score={best_f1:.4f}")
    print('-------------------------------------')


    print("[*] Testing model with the best threshold from validation found...")
    results = test_model(test_packets, bloom, weights, n, best_t) 

    # save results to a file
    with open(f'test_results_{n}.txt', 'w') as f:
        f.write("Index\tScore\tLabel\n")
        for index, score, label in results:
            f.write(f"{index}\t{score:.6f}\t{label}\n")
    print(f"[*] Test results saved to test_results_{n}.txt")

   
if __name__ == "__main__":
    sheet3_task1()
\end{lstlisting}